\scanpage{61}
\margin{\rotatebox{180}{60}}

$F_0$ le Module qu'il induit sur $X_0$. On suppose a) Pour tout $s\quad S$,
la section $t_s$ induite par $t$ sur la fibre $X_s$ est
$\underline O_{X_s}$-régulière et $\underline F_s$-régulière
b) $\underline F$ et $\underline O_X$ sont plats par rapport à $S$
c) Pour tout $s\quad S$, $F_{0\,s}$, \deleted{et}
$\underline O_{X_{0\,s}}$ \inserted{et $\underline O_X$}
sont de profondeur $\quad 2$ en les points fermés de
$X_s\deleted{\unsure{s|x}}$.
Sous ces conditions, il existe un Module cohérent $F$ sur $X$, et un isomorphisme
de son complété formel $\hat F$ le long de $X_0$ avec $\underline F$.

Cet énoncé va résulter du suivant:

\uline{Corollaire} 3.2.\quad Sous les conditions de \unsure{3;1|3.1}.\quad (

\[
\sum_{p\geq0}R^i\unsure{f|g}_{0*}F_0(q-p)
\qquad q-p\geq\unsure{m|n}
\]
\[
S_{\unsure{M|pq}}=
\begin{cases}0\text{ si }p>0\\S_q\text{ si }p=0\end{cases}
\]

\figure{\begin{tikzpicture}[x=.8cm,y=.8cm]
\draw[thin] (0,3.8) .. controls (0,2.8) and (.2,2) .. (.05,-1.1);
\draw[thin] (0,.8) -- (5.8,.85);
\draw[thin,<-] (-.8,0) .. controls (1,.08) and (4,-.1) .. (6.2,-.08);
\draw (0,0) .. controls (1.6,-.15) and (4.8,-.05) .. (6.1,-.16)
 -- (6.15,.15) .. controls (6.2,-.25) and (6.15,-.3) .. (5.9,-.3)
 .. controls (3.3,-.02) and (1.1,-.1) .. (0,-.25) -- cycle;
\node at (-.65,4) {$q$};
\node at (6.4,-.1) {$p$};
\node at (5,3.45) {$F_0(-p)$};
\begin{scope}[shift={(9.2,-.2)}]
\draw[thin] (0,4) .. controls (.15,2.4) and (-.05,2) .. (.05,-1.1)
 .. controls (.2,-1.2) and (.1,-2.3) .. (0,-2.8);
\draw[thin] (-.95,.1) .. controls (0,-.05) and (2.5,.15) .. (5.3,0);
\draw[thin] (-1.1,-3.7) -- (6.2,2.05);
\end{scope}
\end{tikzpicture}}
